\section{Теоретические основы}

\subsection{PTX как промежуточное представление CUDA}

\subsubsection{Место PTX в стеке компиляции CUDA}

Компиляция CUDA-программы проходит через многоуровневый конвейер.
На входе компилятор \texttt{nvcc}~\cite{NVCC} принимает файл \texttt{.cu},
содержащий смешанный код: C++-код хоста и тела ядер устройства.
Внешний интерфейс компилятора выделяет код устройства и передаёт его в NVVM~--- LLVM-based IR~\cite{Lattner2004}
для GPU~--- который транслируется в PTX.

PTX (Parallel Thread Execution)~--- аппаратно-независимое промежуточное
представление, описывающее программу в терминах виртуального RISC-процессора
с бесконечным регистровым файлом и явной параллельностью потоков.
На следующем шаге \texttt{ptxas} транслирует PTX в SASS~---
бинарный машинный код конкретной архитектуры (Volta, Turing, Ampere,
Hopper), который помещается в fatbinary-архив вместе с исходным PTX-текстом.
При запуске CUDA Runtime JIT-компилирует PTX для установленного GPU
при отсутствии подходящего SASS.

Fatbinary является ELF-совместимым архивом, встраиваемым компоновщиком
в секцию \texttt{.nv\_fatbin} результирующего ELF-файла приложения.
Эмулятор извлекает PTX-текст из этой секции при перехвате
\texttt{\_\_cudaRegisterFatBinary}, вызывая утилиту
\texttt{cuobjdump --dump-ptx} (подробно~--- в Главе~3).

\begin{figure}[h]
  \centering
  \includegraphics[width=16cm]{cuda_stackbw.png}
  \caption{Стек компиляции CUDA. PTX занимает уровень между платформно-независимым IR и машинным кодом SASS.}
  \label{fig:cuda_stack}
\end{figure}


\subsubsection{Cистема типов PTX}

PTX определяет следующие базовые типы данных. Тип операнда определяется
двумя способами: суффиксом инструкции (\texttt{.f32}, \texttt{.s32} и др.)
и префиксом имени регистра (\texttt{\%f}~$\to$~\texttt{f32},
\texttt{\%r}~$\to$~\texttt{s32}, \texttt{\%rd}~$\to$~\texttt{u64},
\texttt{\%p}~$\to$~\texttt{pred}):

\begin{itemize}
  \item \textbf{Целочисленные знаковые:} \texttt{s8}, \texttt{s16},
        \texttt{s32}, \texttt{s64}~--- арифметика в дополнительном коде.
  \item \textbf{Целочисленные беззнаковые:} \texttt{u8}, \texttt{u16},
        \texttt{u32}, \texttt{u64}.
  \item \textbf{Битовые:} \texttt{b8}, \texttt{b16}, \texttt{b32},
        \texttt{b64}, \texttt{b128}~--- интерпретация не задана,
        используются для пересылок и упаковки.
  \item \textbf{Вещественные:} \texttt{f16}, \texttt{f32}, \texttt{f64}~---
        форматы IEEE 754-2008~\cite{IEEE754-2008} половинной, одинарной и двойной точности.
  \item \textbf{Предикат:} \texttt{pred}~--- однобитное логическое значение,
        хранимое в регистрах \texttt{\%p0}--\texttt{\%pN}.
\end{itemize}

Суффикс типа однозначно задаёт битовую ширину $w \in \{8, 16, 32, 64, 128\}$
и способ интерпретации битового шаблона.

Различие знаковых, беззнаковых и битовых вариантов одной ширины
непосредственно влияет на реализацию исполнителей: \texttt{shr.u32}
выполняет логический сдвиг, тогда как \texttt{shr.s32}~---
арифметический (с распространением знакового бита).

\subsubsection{Классификация инструкций PTX}

Для целей проектирования таблицы диспетчеризации все инструкции PTX
разбиты на семь классов:

\begin{enumerate}
  \item \textbf{Арифметические} (\texttt{add}, \texttt{sub}, \texttt{mul},
        \texttt{div}, \texttt{fma}, \texttt{mad}, \texttt{rem}, \texttt{abs},
        \texttt{neg}, \texttt{min}, \texttt{max})~--- 25~мнемоник.
  \item \textbf{Логические и сдвиги} (\texttt{and}, \texttt{or},
        \texttt{xor}, \texttt{not}, \texttt{shl}, \texttt{shr},
        \texttt{bfe}, \texttt{bfi}, \texttt{popc})~--- 14~мнемоник.
  \item \textbf{Сравнения и предикаты} (\texttt{setp}, \texttt{selp},
        \texttt{slct}, \texttt{vote})~--- 4~мнемоники.
  \item \textbf{Работа с памятью} (\texttt{ld}, \texttt{st},
        \texttt{mov}, \texttt{cvta}, \texttt{prefetch})~--- 5~мнемоник.
  \item \textbf{Преобразование типов} (\texttt{cvt})~--- 1~мнемоника
        с матрицей суффиксов $5 \times 5$.
  \item \textbf{Управление потоком} (\texttt{bra}, \texttt{call},
        \texttt{ret}, \texttt{exit}, \texttt{trap})~--- 5~мнемоник.
  \item \textbf{Синхронизация и специальные} (\texttt{bar}, \texttt{membar},
        \texttt{mov.special}, \texttt{ex2}, \texttt{lg2},
        \texttt{rcp}, \texttt{sqrt}, \texttt{sin}, \texttt{cos})~--- 9~мнемоник.
\end{enumerate}

\subsubsection{Директивы PTX и их роль в эмуляторе}

В PTX-тексте директивы не являются исполняемыми инструкциями,
но несут информацию, критически необходимую для анализатора:

\begin{itemize}
  \item \texttt{.entry} / \texttt{.func}~--- задают тип функции
        (ядро устройства или вспомогательная функция) и её границы.
  \item \texttt{.param}~--- описывают сигнатуру: имя, тип и выравнивание
        каждого параметра ядра.
  \item \texttt{.reg .type \%prefix<N>}~--- объявляют регистровый файл:
        например, \texttt{.reg .f32 \%f<115>} резервирует 115~регистров
        типа \texttt{f32} с именами \texttt{\%f0}--\texttt{\%f114}.
  \item \texttt{.shared .align A .b8 name[S]}~--- резервирует $S$~байт
        разделяемой памяти с выравниванием $A$~байт.
  \item \texttt{.file ID "path"} и \texttt{.loc ID L C}~--- формируют
        таблицу, привязывающую PTX-инструкции к строкам исходного кода,
        используемую профилировщиком для аннотации отчёта.
\end{itemize}

\subsubsection{Пространства адресов и их семантика}

PTX определяет шесть пространств адресов, отличающихся по видимости
и времени жизни:

\begin{itemize}
  \item \texttt{.reg}~--- регистровый файл, приватный для каждого потока.
  \item \texttt{.global}~--- глобальная память, разделяемая всеми блоками.
  \item \texttt{.shared}~--- разделяемая память, видимая всем потокам
        одного блока; время жизни~--- время жизни блока.
  \item \texttt{.local}~--- приватная память стека потока (спилл регистров).
  \item \texttt{.param}~--- область передачи параметров ядра с хоста.
  \item \texttt{.const}~--- константная память, только для чтения.
\end{itemize}

Формально вводится функция пространства
\[
  \mathit{space} : \mathit{Addr} \to
    \{\mathit{global},\;\mathit{shared},\;\mathit{local},\;\mathit{param},\;\mathit{const}\},
\]
используемая исполнителями \texttt{ld} и \texttt{st} для выбора
соответствующего обработчика обращения к памяти.

Смешение пространств в реализации исполнителей недопустимо:
глобальная и разделяемая память обслуживаются разными аппаратными блоками
с различными правилами доступа, что непосредственно влияет на вычисление
метрик слияния транзакций и конфликтов банков.

\newpage
\subsection{Математическая модель SIMT-исполнения}

\subsubsection{Варп как алгебраическая структура}

В модели SIMT все потоки варпа в каждый момент времени исполняют
одну и ту же инструкцию, однако каждый поток имеет \emph{собственный
счётчик команд и регистровый файл}. Это допускает дивергенцию:
при условном переходе потоки, для которых условие выполнено, и потоки,
для которых оно не выполнено, временно расходятся по разным ветвям,
а аппаратура отслеживает активность каждого потока через маску исполнения
и стек реконвергенции.

Базовой единицей исполнения на GPU является \textit{варп}~---
группа из 32 потоков, всегда исполняющих одну и ту же инструкцию~\cite{Lindholm2008, Gergel2010}.
Формально варп определяется как кортеж:
\begin{equation}
  W = \langle T,\, M,\, \mathrm{PC},\, R,\, S \rangle,
\end{equation}
где:
\begin{itemize}
  \item $T = \{0, 1, \ldots, 31\}$~--- множество идентификаторов потоков варпа;
  \item $M \in \mathbb{B}^{32}$~--- маска исполнения (execution mask):
        $M[i] = 1$ означает, что поток $i$ исполняет целевую инструкцию;
  \item $\mathrm{PC} \in \mathbb{N}$~--- программный счётчик;
  \item $R$~--- регистровый файл (функция $T \times \mathit{Types} \times \mathbb{N} \to \mathit{Val}$);
  \item $S$~--- стек дивергенции (подробно в Главе~6).
\end{itemize}

Операция исполнения одной инструкции определяется как функция перехода
состояния:
$\delta : W \times \mathit{Instr} \to W'$.

\subsubsection{Маска исполнения: операции и инварианты}

Маска $M \in \mathbb{B}^{32}$ формально задаёт, какие потоки исполняют
следующую инструкцию. Определим вспомогательные операции:
\begin{align}
  \activeset(M) &= \{i \in T \mid M[i] = 1\}, \\
  \popcount(M) &= |\activeset(M)|, \\
  \mathit{apply}(M, M') &= M \mathbin{\wedge} M'
    \quad\text{(маскирование по предикату)}.
\end{align}

Фундаментальный инвариант исполнения: \textit{исполнитель инструкции
обязан игнорировать потоки с $M[i]=0$ при записи результата.}
Нарушение данного инварианта порождает логические ошибки
в реализации исполнителей и строго контролируется в тестах.

В отличие от наивного однопоточного исполнителя, потоки с $M[i]=0$
не просто пропускают инструкцию~--- они не должны изменять регистровый
файл и память, иначе результат разойдётся с аппаратным SIMT-исполнением.

\subsubsection{Предикатные регистры и составная маска}

Предикатный регистр $P \in \mathbb{B}^{32}$ хранит по одному биту
на каждый поток варпа. Инструкция \texttt{@\%p instr} выполняется
с эффективной маской:
\begin{equation}
  M_{\mathrm{eff}} = M \mathbin{\wedge} P,
\end{equation}
а инструкция \texttt{@!\%p instr}~--- с:
\begin{equation}
  M_{\mathrm{eff}} = M \mathbin{\wedge} \neg P.
\end{equation}

Инструкция \texttt{setp.op.type \%p, a, b} вычисляет:
\begin{equation}
  P[i] = \mathit{op}(a_i, b_i), \quad \forall i \in \activeset(M),
\end{equation}
что при следующем условном переходе \texttt{@\%p bra} может привести
к расщеплению варпа, если $M \wedge P \neq 0$ и $M \wedge \neg P \neq 0$.

\subsection{Поведение инструкций PTX}

\subsubsection{Обозначения}

Состояние варпа описывается тремя компонентами:
\begin{itemize}
  \item $R[t][\tau][k]$~--- значение $k$-го регистра типа $\tau$ потока $t$;
  \item $M$~--- 32-битная маска исполнения (бит~$t$ равен~1, если поток~$t$ активен);
  \item $\mathrm{PC}$~--- номер текущей инструкции.
\end{itemize}

Запись $\forall t \in \activeset(M)$ означает, что правило применяется
только к активным потокам; регистры неактивных потоков не изменяются.

\subsubsection{Арифметические инструкции}

\texttt{add.s32 \%rd, \%ra, \%rb} складывает два 32-битных целых
и записывает результат в $\mathtt{\%rd}$, отбрасывая переполнение:
\begin{equation}
  R'[t][\mathtt{s32}][d] =
  \bigl(R[t][\mathtt{s32}][a] + R[t][\mathtt{s32}][b]\bigr) \bmod 2^{32},
  \quad \forall t \in \activeset(M).
\end{equation}

\texttt{fma.rn.f32 \%fd, \%fa, \%fb, \%fc} вычисляет $a \cdot b + c$
за одну операцию с одним округлением по стандарту IEEE~754~\cite{IEEE754-2008}
(режим «к ближайшему чётному»):
\begin{equation}
  R'[t][\mathtt{f32}][d] =
  \mathrm{round}_{n}\!\bigl(R[t][\mathtt{f32}][a]
    \cdot R[t][\mathtt{f32}][b]
    + R[t][\mathtt{f32}][c]\bigr),
  \quad \forall t \in \activeset(M).
\end{equation}

\subsubsection{Инструкции сравнения и выбора}

\texttt{setp.lt.s32 \%p, \%ra, \%rb} записывает результат сравнения
в предикатный регистр $\mathtt{\%p}$:
\begin{equation}
  P'[t] = \bigl(R[t][\mathtt{s32}][a] < R[t][\mathtt{s32}][b]\bigr),
  \quad \forall t \in \activeset(M).
\end{equation}

\texttt{selp.f32 \%fd, \%fa, \%fb, \%p} выбирает одно из двух значений
в зависимости от предиката, не вызывая дивергенцию:
\begin{equation}
  R'[t][\mathtt{f32}][d] =
  \begin{cases}
    R[t][\mathtt{f32}][a] & \text{если } P[t] = 1, \\
    R[t][\mathtt{f32}][b] & \text{если } P[t] = 0,
  \end{cases}
  \quad \forall t \in \activeset(M).
\end{equation}

\subsubsection{Инструкции управления потоком}

\texttt{bra.uni \$L} — безусловный переход: все потоки варпа
переходят к метке \$L, маска~$M$ не меняется:
\begin{equation}
  \mathrm{PC}' = \mathit{offset}(\$L),
\end{equation}
где $\mathit{offset}(\$L)$~--- номер первой инструкции метки \$L
из таблицы базовых блоков.

\texttt{@\%p bra \$L} — условный переход. Обозначим $M_{\mathrm{taken}} = M \wedge P$
(потоки, берущие переход) и $M_{\mathrm{not}} = M \wedge \neg P$ (остальные):

\begin{itemize}
  \item Если $M_{\mathrm{not}} = 0$: все активные потоки берут переход,
        $\mathrm{PC}' = \mathit{offset}(\$L)$.
  \item Если $M_{\mathrm{taken}} = 0$: ни один поток не берёт переход,
        $\mathrm{PC}' = \mathrm{PC} + 1$.
  \item Иначе: часть потоков расходится — применяется алгоритм
        PDOM\_Diverge (подробно в Главе~6).
\end{itemize}

\subsubsection{Инструкции памяти}

\texttt{ld.global.f32 \%fd, [\%ra]} читает 4-байтное значение
из глобальной памяти по адресу в регистре \texttt{\%ra}:
\begin{equation}
  R'[t][\mathtt{f32}][d] = \mathit{Mem}_{\mathrm{global}}[R[t][\mathtt{u64}][a]],
  \quad \forall t \in \activeset(M).
\end{equation}

После исполнения \texttt{ld.global}/\texttt{st.global} вычисляются
$T_{\mathrm{actual}}$ и $\kappa$ по формулам~\eqref{eq:t_actual_ch7}--\eqref{eq:coalescing_ch7}.

\subsection{Механизм динамической подстановки библиотек в Linux}

\subsubsection{Загрузчик ELF и таблица динамических символов}

При запуске ELF-бинарного файла в Linux ядро передаёт управление
динамическому загрузчику \texttt{ld.so}. Загрузчик строит
\textit{список разрешения символов}~--- упорядоченную цепочку библиотек,
в которой он ищет адрес каждого используемого символа~\cite{LinuxManLdso, Stolyarov2011}.

При помещении \texttt{libemuruntime.so} на позицию $l_1$
(через \texttt{LD\_PRELOAD}) вызовы CUDA Runtime~---
\texttt{cudaMalloc}, \texttt{cudaLaunchKernel} и прочие~---
разрешаются в эмуляторе по правилу первого совпадения,
не достигая \texttt{libcudart.so}.


\subsection{Обзор смежных работ}

\subsubsection{Функциональные эмуляторы: Ocelot, CUDAsim}

Проект Ocelot~\cite{Diamos2010} разработан в Технологическом институте
Джорджии в 2009--2013~гг. Он реализует функциональную эмуляцию PTX
с поддержкой перехвата CUDA~4.x API через механизм, аналогичный
\texttt{LD\_PRELOAD}. Ocelot не поддерживает PTX~ISA~7.x+ (инструкции
\texttt{ldmatrix}, \texttt{wmma}, \texttt{async.bulk.copy})
и не компилируется с современными версиями GCC/Clang. Возможности
профилирования ограничены счётчиком выполненных инструкций.

CUDAsim (CUDA Device Emulation, устаревший режим)~--- встроенный в CUDA~4.x
режим эмуляции CPU, удалённый в CUDA~5.0 из-за расхождений в поведении
атомарных операций и барьеров.

\subsubsection{Микроархитектурный симулятор GPGPU-Sim}

GPGPU-Sim~\cite{Bakhoda2009}~--- timing-accurate симулятор, воспроизводящий
конвейер SM с точностью до цикла: внешний интерфейс, исполнение инструкций,
исполнительные блоки, кэши, DRAM. Поддерживаемое подмножество PTX~ISA
ограничено версиями до~4.x; инструкции PTX~7.x+
(\texttt{ldmatrix}, \texttt{wmma}, атрибут \texttt{.address\_size~64})
не поддерживаются. Его цель принципиально отличается от задачи данной работы:
оценка латентностей и пропускной способности, а не функциональная
корректность в контексте непрерывной интеграции. Один запуск \texttt{mmul\_custom} (64×64×64)
в GPGPU-Sim занимает порядка 10~часов машинного времени.

\textit{Сравнение подходов} приведено в таблице~\ref{tab:related_work}.

\begin{table}[h]
  \caption{Сравнение инструментов эмуляции/симуляции CUDA}
  \label{tab:related_work}
  \centering
  \footnotesize
  \begin{tabularx}{\linewidth}{p{2cm}XXXX}
    \toprule
    \textbf{Инструмент} & \textbf{PTX ISA} & \textbf{Перехват} &
    \textbf{Профилировщик} & \textbf{Тип} \\
    \midrule
    Ocelot           & 1.x--2.x  & LD\_PRELOAD  & Счётчики & Функциональный \\
    GPGPU-Sim        & $\leq$~4.x & Не требуется & Timing   & Архитектурный \\
    Compute Sanitizer & —         & Требует GPU  & Нет      & Аппаратный \\
    \textbf{Данная работа} & \textbf{7.x} & \textbf{LD\_PRELOAD} &
    \textbf{4 метрики} & \textbf{Функциональный} \\
    \bottomrule
  \end{tabularx}
\end{table}

\newpage
